Put
\[
 v_-:=\min_{d\in\mathcal D}\min\{V_{E,d},V_{O,d}\}>0,
 \qquad
 v_+:=\max_{d\in\mathcal D}\max\{V_{E,d},V_{O,d}\}<\infty,
\]
and
\[
 Z_m:=\max_{d\in\mathcal D}\max\{|\widehat V_{E,d}-V_{E,d}|,
 |\widehat V_{O,d}-V_{O,d}|\}.
\]
For \(x,y>0\), define
\[
 F_d(x,y)=\left(\sqrt{c_{E,d}x}+\sqrt{c_{O,d}y}\right)^2.
\]
The budget-profile theorem gives \(F_d(V_{E,d},V_{O,d})=\mathcal V_d^*(P)\), while the pilot construction gives \(F_d(\widehat V_{E,d},\widehat V_{O,d})=\widehat{\mathcal V}_{d,m}^*\).  Since \(\mathcal D\) is finite, the gradients of all \(F_d\) are uniformly bounded on \([v_-/2,2v_+]^2\).  On \(\{Z_m\le v_-/2\}\), the mean-value theorem therefore gives \(\varepsilon_m\le C Z_m\) for a finite constant \(C\).  The pilot variance assumptions imply
\[
 Z_m=O_P(r_{V,m})=o_P(1),\qquad
 \varepsilon_m=O_P(r_{V,m})=o_P(1). \tag{1}
\]

Choose any \(d^*\in D^*(P)\).  Empirical optimality of \(\widehat d_m\) gives, pathwise,
\[
\begin{aligned}
0&\le \mathcal V_{\widehat d_m}^*(P)-\mathcal V_{d^*}^*(P)\\
&\le |\mathcal V_{\widehat d_m}^*(P)-\widehat{\mathcal V}_{\widehat d_m,m}^*|
 +\{\widehat{\mathcal V}_{\widehat d_m,m}^*-\widehat{\mathcal V}_{d^*,m}^*\}
 +|\widehat{\mathcal V}_{d^*,m}^*-\mathcal V_{d^*}^*(P)|\\
&\le2\varepsilon_m. 
\end{aligned} \tag{2}
\]
If \(D^*(P)\ne\mathcal D\), finiteness gives \(\gamma(P)>0\), and misselection makes the left side of (2) at least \(\gamma(P)\).  Hence
\[
 \{\widehat d_m\notin D^*(P)\}\subseteq
 \{2\varepsilon_m\ge\gamma(P)\}.
\]
Equation (1) proves the asserted probability convergence.  If \(D^*(P)=\mathcal D\), both the misselection event and, under the convention \(\gamma(P)=+\infty\), the event on the right are empty.

For the allocation map, define
\[
 A_d(x,y)=\frac{\sqrt{c_{E,d}x}}
 {\sqrt{c_{E,d}x}+\sqrt{c_{O,d}y}}.
\]
The elementary inequality
\[
 |(x\vee\underline v_m)-V|\le |x-V|+\underline v_m
\]
and a second finite-family mean-value argument on the same positive rectangle imply
\[
 \max_{d\in\mathcal D}|\widehat\alpha_{d,m}-\alpha_d^*(P)|
 =O_P(r_{V,m}+\underline v_m)=o_P(1). \tag{3}
\]
All population allocations are uniformly interior because the family is finite and both variance components and costs are strictly positive and finite.  Thus, with probability tending to one, every \(\widehat\alpha_{d,m}\) belongs to one common interval \([a_0,1-a_0]\), where \(a_0>0\).

Uniformly over this interval and \(d\in\mathcal D\), the floor inequality from the budget-profile proof gives
\[
\begin{aligned}
0\le{}&B\left\{\frac{V_{E,d}}{\lfloor\beta B/c_{E,d}\rfloor}
 +\frac{V_{O,d}}{\lfloor(1-\beta)B/c_{O,d}\rfloor}\right\}
 -\mathcal V_d(\beta;P)\\
&\le\frac2B\left\{\frac{V_{E,d}c_{E,d}^2}{\beta^2}
 +\frac{V_{O,d}c_{O,d}^2}{(1-\beta)^2}\right\}=O(B^{-1}).
\end{aligned} \tag{4}
\]
The square identity in the budget-profile proof, or a second-order Taylor expansion at its unique interior minimizer, also gives uniformly in \(d\),
\[
 \mathcal V_d(\widehat\alpha_{d,m};P)-\mathcal V_d^*(P)
 =O_P\!\left((r_{V,m}+\underline v_m)^2\right)
 =O_P(r_{V,m}+\underline v_m). \tag{5}
\]
Combining (4) and (5), and taking \(\beta=\widehat\alpha_{d,m}\), proves the asserted uniform
\[
 O_P(r_{V,m}+\underline v_m+B^{-1})=o_P(1)
\]
bound.  The same interior event makes all displayed floor counts positive for sufficiently large \(B\).  The deterministic tie rule affects only which empirical minimizer is implemented; equation (2) and the separation argument concern membership in the entire population optimizer set.